\section*{Introduction}

While Martin-L\"of Type Theory (MLTT) is a synthetic language for $\infty$-groupoids,
\emph{directed type theory} -- which is by no means a fixed construct -- 
aims to achieve the same for higher categories.

\subsection*{Size issues}

Let us assume that we have three sizes: small, large, and extra-large.
In classical foundations, this could be done by postulating enough
strongly inaccessible cardinals or
(equivalently) Grothendieck universes \cite{bourbaki1972}.
We denote the 1-category of small categories as $\cat$,
which is an object in the 1-category of large categories as $\Cat$.
We will also use $\tCat$, which differs from the others in two ways:
$\tCat$ is a strict 2-category rather than a 1-category,
and its objects are extra-large categories.

\subsection*{Overview of the literature}

There is an analogue of the Hofmann--Streicher groupoid model of MLTT \cite{hofmann1995},
where the category of categories $\Cat$ replaces the category of groupoids
as the category of contexts.
% There is a 1-categorical approach to the semantics of directed type theory,
% emulating Hofmann and Streicher's groupoid model of MLTT \cite{hofmann1995}.
North \cite{north2019} proposed a type theory for this model,
with \emph{hom-types} replacing identity types,
placing restrictions on
the introduction and elimination rules for hom-types using the groupoid core of a category.
% The model in $\Cat$ then makes use of a certain
% algebraic weak factorisation system on $\Cat$.
Altenkirch and Neumann \cite{altenkirch2024} continued this work by
formulating directed type theory in terms of a generalised algebraic theory,
providing several category-theoretic constructions in this synthetic category theory.
Chu and North \cite{rivera2026} provided an analysis of dependent 2-sided fibrations,
allowing reasoning with natural transformations,
which proved challenging in \cite{altenkirch2024}.
Also notable is their treatment of \emph{hom-types},
which no longer placed restrictions on introduction and elimination rules
using the groupoid core.

For higher-categorical models,
the model structure for complete Segal spaces in bisimplicial sets
presents the category theory of $\infty$-categories \cite{riehlverity2017}.
Riehl and Shulman \cite{riehl2017} developed a three-layered type theory,
with \emph{Simplicial Type Theory} at the highest level,
and an interpretation in complete Segal spaces.
Not all types in Simplicial Type Theory model higher categories.
Rather, $\infty$-categories are modelled by Rezk types (and Segal types).
There is a growing list of results developed in this setting,
which we do not detail here
(overviews are provided in \cite{riehl2025, gratzer2026}, for example).
Some of these results have also been formalised in the proof assistant Rzk \cite{rzk}.

A bicategorical approach to directed type theory is also natural to consider.
Licata and Harper \cite{licata2011} defined 2-Dimensional Directed Type Theory
that has a syntactic notion of a 2-cell.
The authors developed a syntax for directed $\Pi$-types
that can handle the complexity introduced by the variances (polarity) of types.
An ad-hoc model of 2-Dimensional Directed Type Theory was
then constructed in $\tCat$.
Ahrens et al. \cite{ahrens2023} defined \emph{comprehension bicategories},
capturing what it means to be a general model of 2-Dimensional Directed Type Theory,
for which the $\tCat$ is an example.

Laretto et al. \cite{laretto2026} defined a non-dependent directed type theory,
where types are modelled by categories,
special ``propositional contexts'' are modelled by profunctors,
and their associated substitutions are modelled by dinatural transformations.
Notably, this work includes an axiomatisation of ends and coends.

Our work is based upon the MLTT-style approach of \cite{north2019, altenkirch2024},
with the same focus on the model in the 1-category $\Cat$.
We will adopt the methods of \emph{algebraic type theory} \cite{awodey2025}
to analyse type formers in this setting,
using the polynomial machinery developed in \Cref{sec:polynomials}.
The aim here is not necessarily to prove something new about the model,
but rather to reformulate existing results into a more algebraic language
that can illuminate what is needed for a full axiomatisation.

\subsection*{Challenges and limitations}

An axiomatisation of the model in $\Cat$ is given in \Cref{sec:type-theoretic-axioms},
and it is by no means complete.
For example, the axioms fail to capture the fact that
the formation of $\Sigma$-types and $\Pi$-types
preserves morphisms of opfibrations.
This is discussed further in \Cref{rmk:pi-preserve-opfib-mor}.
There is no treatment of univalence here,
which remains a challenge.
There is also no consideration of directed higher inductive types.
The theory of $W$-types in $\Cat$
would be an interesting application of the theory of polynomials in $\Cat$
that we develop.

Before jumping into the model in $\Cat$,
we make a summary of the various intricacies of the model,
which will be explained in further detail during the exposition.
\begin{itemize}
  \item There are two kinds of types,
    corresponding to two kinds of variances of indexed categories,
    or two kinds of fibrations.
  \item Whilst there are the usual $\Unit$ and $\Sigma$-types for both variances,
    the $\Pi$-types in $\Cat$ involve an interaction between the two variances.
  \item $\Hom$-types give rise to twisted arrow categories,
    rather than simply arrow categories,
    and this results in restricted introduction and elimination rules
    associated with $\Hom$-types.
  \item Like in the category of groupoids,
    opfibrancy and fibrancy can be viewed as types having transport.
    Unlike in the category of groupoids,
    not all morphisms in a slice between opfibrations preserve transport,
    which is to say that not all morphisms (in a slice) between opfibrations
    are morphisms of opfibrations.
    We do not address this issue in this work.
\end{itemize}

\section{A case study in the category of categories}
\label{sec:preliminaries}
We recall some of the basic definitions in $\Cat$
that will be relevant to our model.
There are two interacting classes of maps,
namely the fibrations and opfibrations.
We will focus on opfibrations,
and only introduce notation for the dual notion of
fibrations when needed.

\section*{Opfibrations}
% \label{sec:fibrations-in-cat}

Let $\DD$ and $\CC$ be categories and
$F : \DD \to \CC$ be a functor.
We say $f : x \to y$ is a lift of its image $F f : F x \to F y$.
Recall that $f'$ is \emph{opcartesian}
when for every morphism $g : y \to z$ in $\CC$
and every lift $h' : x' \to z'$ of $g \circ f : x \to z$
there is a unique lift $g' : y' \to z'$ of $g$ such that
$h' = g' \circ f'$.
\[ \begin{tikzcd}[row sep = tiny]
  && {z'} \\
  {x'} & {y'} &&& \DD \\
  && z \\
  x & y &&& \CC
  \arrow["{h'}", from=2-1, to=1-3]
  \arrow["{f'}"', from=2-1, to=2-2]
  \arrow[dashed, from=2-2, to=1-3]
  \arrow["F", from=2-5, to=4-5]
  \arrow[from=4-1, to=3-3]
  \arrow["f"', from=4-1, to=4-2]
  \arrow["g"', from=4-2, to=3-3]
\end{tikzcd}\]
% \end{definition}

Recall also the following ``pullback-hom'' diagram in $\Cat$,
where the pullback computes a comma category $F / \CC$.
\begin{equation}\label{eq:pullback-hom}
\begin{tikzcd}
  {\DD^{\two}} \\
  & {F / \CC} & \DD \\
  & {\CC^{\two}} & \CC
  \arrow["{h_F}"{description}, dashed, from=1-1, to=2-2]
  \arrow["\src", bend left, from=1-1, to=2-3]
  \arrow["{F^{\two}}"', bend right, from=1-1, to=3-2]
  \arrow[from=2-2, to=2-3]
  \arrow[from=2-2, to=3-2]
  \arrow["\lrcorner"{anchor=center, pos=0.125}, draw=none, from=2-2, to=3-3]
  \arrow["F", from=2-3, to=3-3]
  \arrow["\src"', from=3-2, to=3-3]
\end{tikzcd}
\end{equation}

The functor $h_F$ takes an arrow $f : x \to y$ in $\DD$
and constructs a lifting problem (i.e. an element of the comma category)
by taking the domain of the arrow in $\DD$ and the image of the arrow in $\CC$.
\[ h_F (f) = (x,F f)\]
% The same kind of diagram can be set up using the twisted arrow category
% \Cref{def:twisted-arrow} instead,
% to construct a ``twisted pullback-hom'' into a ``twisted comma category''
% $F \tcomma \CC$.
% \[ \begin{tikzcd}
%     {\Tw{\DD}} && \\
%     & {F \tcomma \CC} & {{\DD}^\op} \\
%     & {\Tw{\CC}} & {{\CC}^\op}
%     \arrow["{{t_F}}"{description}, dashed, from=1-1, to=2-2]
%     \arrow["\src", bend left, from=1-1, to=2-3]
%     \arrow["{\Tw{F}}"', bend right, from=1-1, to=3-2]
%     \arrow[from=2-2, to=2-3]
%     \arrow[from=2-2, to=3-2]
%     \arrow["\lrcorner"{anchor=center, pos=0.125}, draw=none, from=2-2, to=3-3]
%     \arrow["F", from=2-3, to=3-3]
%     \arrow["\src"', from=3-2, to=3-3]
%   \end{tikzcd}\]

% \begin{proposition}\label{prop:hurewicz-equiv-defs}
%   Let $F : \DD \to \CC$ be a functor.
%   The following are equivalent (structures)
%   \begin{itemize}
%     \item For any functors
%       $p : {\two} \to \CC$
%       and $x : 1 \to \DD$
%       such that $F \circ x = p \circ \delta_0$,
%       a chosen lift $l_{(x,p)} : {\two} \to \DD$
%       satisfying $l \circ \delta_0 = x$ and $F \circ l = p$.
%       \[ \begin{tikzcd}
%         {1} & \DD \\
%         {{\two}} & \CC
%         \arrow["x", from=1-1, to=1-2]
%         \arrow["\delta_0"', from=1-1, to=2-1]
%         \arrow["F", from=1-2, to=2-2]
%         \arrow["l"{description}, dashed, from=2-1, to=1-2]
%         \arrow["p"', from=2-1, to=2-2]
%       \end{tikzcd} \]
%     \item For each category $\EE$, functors
%       $p : {\two \times \EE} \to \CC$
%       and $x : \EE \to \DD$
%       such that $F \circ x = p \circ \delta_0$,
%       a chosen lift $l_{(x,p)} : {\two \times \EE} \to \DD$
%       satisfying $l \circ \delta_0 = x$ and $F \circ l = p$.
%       \[ \begin{tikzcd}
%         {\EE} & \DD \\
%         {{\two \times \EE}} & \CC
%         \arrow["x", from=1-1, to=1-2]
%         \arrow["\delta_0"', from=1-1, to=2-1]
%         \arrow["F", from=1-2, to=2-2]
%         \arrow["l"{description}, dashed, from=2-1, to=1-2]
%         \arrow["p"', from=2-1, to=2-2]
%       \end{tikzcd} \]
%       Furthermore for any $G : \EE' \to \EE$,
%       \[l_{(x,p)} \circ (\two \times G) = l_{(x \circ G, p \circ (\two \times G))}\]
%       \[ \begin{tikzcd}
%         {\EE'} & {\EE} & \DD \\
%         {\two \times \EE'} & {{\two \times \EE}} & \CC
%         \arrow["{G}", from=1-1, to=1-2]
%         \arrow["{\delta_0}"', from=1-1, to=2-1]
%         \arrow["x", from=1-2, to=1-3]
%         \arrow[from=1-2, to=2-2]
%         \arrow["F", from=1-3, to=2-3]
%         \arrow[dashed, from=2-1, to=1-3]
%         \arrow["{\two \times G}"', from=2-1, to=2-2]
%         \arrow[dashed, from=2-2, to=1-3]
%         \arrow["p"', from=2-2, to=2-3]
%       \end{tikzcd}\]
%     \item A section $l : F / \CC \to \DD^\two$
%     of the pullback-hom $h_F : \DD^\two \to F / \CC$.
%     \[ \begin{tikzcd}
%        & {\DD^\two} \\
%       {F / \CC} & {F / \CC}
%       \arrow["{h_F}", from=1-2, to=2-2]
%       \arrow["l", from=2-1, to=1-2]
%       \arrow[equals, from=2-1, to=2-2]
%     \end{tikzcd} \]
%     \item For each category $\EE$, functors
%       $p : \Cyl{\EE} \to \CC$
%       and $x : \EE^\op \to \DD$
%       such that $F \circ x = p \circ \inl$,
%       a chosen lift $l_{(x,p)} : \Cyl{\EE} \to \DD$
%       satisfying $l \circ \inl = x$ and $F \circ l = p$.
%       \[ \begin{tikzcd}
%         {\EE^\op} & \DD \\
%         {\Cyl{\EE}} & \CC
%         \arrow["x", from=1-1, to=1-2]
%         \arrow["\inl"', from=1-1, to=2-1]
%         \arrow["F", from=1-2, to=2-2]
%         \arrow["l"{description}, dashed, from=2-1, to=1-2]
%         \arrow["p"', from=2-1, to=2-2]
%       \end{tikzcd} \]
%       Furthermore for any $G : \EE' \to \EE$,
%       \[l_{(x,p)} \circ \Cyl{G} = l_{(x \circ G^\op, p \circ \Cyl{G})}\]
%       \[ \begin{tikzcd}
%         {\EE'} & {\EE} & \DD \\
%         {\two \times \EE'} & {{\two \times \EE}} & \CC
%         \arrow["{G^\op}", from=1-1, to=1-2]
%         \arrow["{\delta_0}"', from=1-1, to=2-1]
%         \arrow["x", from=1-2, to=1-3]
%         \arrow[from=1-2, to=2-2]
%         \arrow["F", from=1-3, to=2-3]
%         \arrow[dashed, from=2-1, to=1-3]
%         \arrow["{\Cyl{G}}"', from=2-1, to=2-2]
%         \arrow[dashed, from=2-2, to=1-3]
%         \arrow["p"', from=2-2, to=2-3]
%       \end{tikzcd}\]
%   \end{itemize}
% \end{proposition}

\begin{proposition}\label{prop:opfibration-equiv-defs}
  Let $F : \DD \to \CC$ be a functor.
  The following are equivalent (structures)
  \begin{enumerate}
    \item For each object $x \in \DD$, each object $y \in \CC$,
    and each morphism $f : F x \to y$ in $\CC$,
    a chosen opcartesian lift $l_{(x,f)} : x \to y'$ of $f$.
    \item A section $l : F / \CC \to \DD^\two$
    of the pullback-hom $h_F : \DD^\two \to F / \CC$,
    that factors through the full subcategory
    $\KK \subseteq \DD^\two$ consisting of opcartesian arrows.
    \[ \begin{tikzcd}
      {\KK} & {\DD^\two} \\
      {F / \CC} & {F / \CC}
      \arrow[hook, from=1-1, to=1-2]
      \arrow["{h_F}", from=1-2, to=2-2]
      \arrow[dashed, from=2-1, to=1-1]
      \arrow["l"{description}, from=2-1, to=1-2]
      \arrow[equals, from=2-1, to=2-2]
    \end{tikzcd} \]
    \item A left adjoint $l : F / \CC \to \DD^\two$
    to $h_F : \DD^\two \to F / \CC$,
    with the unit of the adjunction equal to the identity.
    (This makes $l$ a ``left-adjoint-right-inverse'' or ``lari''.)
    \[ \eta : \id_{F / \CC} = h_F \circ l \]
    % \item A section $l : F \tcomma \CC \to \Tw{\DD}$
    % of the twisted pullback-hom $t_F : \Tw{\DD} \to F \tcomma \CC$,
    % that factors through the full subcategory
    % $\TT \subseteq \Tw{\DD}$ consisting of opcartesian twisted arrows.
    % \[ \begin{tikzcd}
    %   {\TT} & {\Tw{\DD}} \\
    %   {F \tcomma \CC} & {F \tcomma \CC}
    %   \arrow[hook, from=1-1, to=1-2]
    %   \arrow["{t_F}", from=1-2, to=2-2]
    %   \arrow[dashed, from=2-1, to=1-1]
    %   \arrow["l"{description}, from=2-1, to=1-2]
    %   \arrow[equals, from=2-1, to=2-2]
    % \end{tikzcd} \]
    % \item A left-adjoint-right-inverse $l : F \tcomma \CC \to \Tw{\DD}$
    % to $t_F : \Tw{\DD} \to F \tcomma \CC$, with identity unit.
    % \[ \eta : \id_{F \tcomma \CC} = t_F \circ l \]
  \end{enumerate}
\end{proposition}
\begin{proof}
  (2) $\to$ (1) is straightforward by
  defining $l_{(x,f)} := l (x,f)$.
  The equation $\id_{F / \CC} = h_F \circ l$ ensures that
  $l$ produces correct lifts to lifting problems supplied by $F / \CC$.
  % Since $h_F \circ l = \id$,
  % it suffices to construct a counit $\epsilon : \id \to l \circ h_F$,
  % and show that $h_F \epsilon = \id_{h_F}$ and $\epsilon l = \id_l$.

  (3) $\to$ (2).
  The counit of the adjunction supplies the unique map out of
  the lift given by $l$,
  ensuring that $l$ produces opcartesian maps.
  
  (1) $\to$ (3).
  Let $l : F / \CC \to \DD^\two$ be given by taking
  $(x, f) \mapsto l_{(x,f)}$ on objects.
  For a morphism $(\psi : x \to x', \phi : y \to y') : (x,f : F x \to y) \to (x',g : F x' \to y')$,
  we define $l (\psi,\phi) : l_{(x,f)} \to l_{(x',f')}$
  by lifting the commuting square for $(\psi,\phi)$
  using the opcartesian property of $l_{(x,f)}$.
  \[\begin{tikzcd}
    Fx & y & x & {y_0} \\
    {Fx'} & {y'} & {x'} & {y_0'}
    \arrow["f", from=1-1, to=1-2]
    \arrow["{F\psi}"', from=1-1, to=2-1]
    \arrow["\phi", from=1-2, to=2-2]
    \arrow["{l_{(x,f)}}", from=1-3, to=1-4]
    \arrow["\psi"', from=1-3, to=2-3]
    \arrow["{\phi_0}", dashed, from=1-4, to=2-4]
    \arrow["{f'}"', from=2-1, to=2-2]
    \arrow["{l_{(x',f')}}"', from=2-3, to=2-4]
  \end{tikzcd}\]
  More precisely:
  \[l (\psi,\phi) := (\psi,\phi_0) : l_{(x,f)} \to l_{(x',f')}\]
\end{proof}
% Condition $2.$ can be viewed as an extension of the ``Hurewicz fibration'' condition
% used for the semantics of path types \cite{awodey2026}.
% Condition $3.$ is just a more concise rephrasing of $2.$
% Conditions $4.$ and $5.$ will replace $2.$ and $3.$ in our
% axiomatisation of directed path types in \Cref{sec:directed-types}.

\begin{definition}\label{def:split-opfibration}
  A cloven opfibration consists of a pair $(F,l)$,
  where $F : \DD \to \CC$ is a functor
  and $l$ is left-adjoint-right-inverse to
  the pullback-hom $h_F : \DD^\two \to F / \CC$
  (or one of the other equivalent structures of
  \Cref{prop:opfibration-equiv-defs}).

  A cloven opfibration $(F : \DD \to \CC,l)$
  is \emph{normal} if 
  for any $x \in \DD$ the lift of the identity
  $\id_{F x} : F x \to F x$ is the identity
  \[l_{(x,\id_{F x})} = \id_x : x \to x\]
  In particular, the endpoint of the lift satisfies $(F x)' = x$.
  We say $(F,l)$ is a \emph{normal opfibration} for short.

  A normal opfibration $(F : \DD \to \CC, l)$ is \emph{split} if
  for any $x \in \DD$, $y$ and $z$ in $\CC$ and
  morphisms $f : F x \to y$ and $g : y \to z$,
  the lift of the composition is the composition of the lifts
  \[l_{(x,g \circ f)} = l_{(x,g)} \circ l_{(y',f)} : x \to z' \]
  In particular, the endpoints $z'$ of the two lifts agree.
  We call $(F,l)$ a \emph{split opfibration} for short.
\end{definition}

We can of course phrase these conditions more diagrammatically;
we do so for the normality condition.
\begin{proposition}\label{prop:normal-opfibration-equiv-defs}
  Let $(F : \DD \to \CC, l)$ be a cloven opfibration.
  The following conditions are equivalent
  \begin{enumerate}
    \item $(F,l)$ is normal.
    \item The constant functor $\Delta : \DD \to \DD^\two$
      factors as $\Delta = l \circ L_F$, where
      $L_F := h_F \circ \Delta$.
      \[ \begin{tikzcd}
        \DD & {\DD^\two} \\
        {\DD^\two} & {F / \CC}
        \arrow["\Delta", from=1-1, to=1-2]
        \arrow["\Delta"', from=1-1, to=2-1]
        \arrow["{L_F}"{description}, dashed, from=1-1, to=2-2]
        \arrow["{{h_F}}"', from=2-1, to=2-2]
        \arrow["l"', from=2-2, to=1-2]
      \end{tikzcd} \]
    % \item The functor $\rfl : \Core{\DD} \to \Tw{\DD}$ that takes
    %   an object $d$ to the identity on $d$ and an isomorphism
    %   $f : d \to d'$ to the twisted morphism
    %   $(f^{-1}, f) : \id_d \to \id_d$
    %   factors as $\rfl = l \circ t_F \circ \rfl$.
    %   \[ \begin{tikzcd}
    %     d & {d'} && {\Core{\DD}} & {\Tw{\DD}} \\
    %     d & {d'} && {\Tw{\DD}} & {F \tcomma \CC}
    %     \arrow["{=}"', from=1-1, to=2-1]
    %     \arrow["{f^{-1}}"', from=1-2, to=1-1]
    %     \arrow["{=}", from=1-2, to=2-2]
    %     \arrow["\rfl", from=1-4, to=1-5]
    %     \arrow["\rfl"', from=1-4, to=2-4]
    %     \arrow["f"', from=2-1, to=2-2]
    %     \arrow["{t_F}"', from=2-4, to=2-5]
    %     \arrow["l"', from=2-5, to=1-5]
    %   \end{tikzcd} \]
  \end{enumerate}
\end{proposition}
\begin{proof}
  (1) $\implies$ (2).
  For an object $x \in \DD$,
  \[l \circ h_F \circ \Delta (x) = l \circ h_F (x,\id_x) = l (x, \id_{F x}) = \id_x = \Delta (x)\]  
  For a morphism $f : x \to y$ in $\DD$,
  \[l \circ h_F \circ \Delta (f) = l \circ h_F (f, f) = l (f, F f) = (f,f) = \Delta (f)\]
  \[ \begin{tikzcd}
    x & x \\
    y & y
    \arrow["{l_{(x,\id)}}", equals, from=1-1, to=1-2]
    \arrow["f"', from=1-1, to=2-1]
    \arrow["f", dashed, from=1-2, to=2-2]
    \arrow["{l_{(y,\id)}}"', equals, from=2-1, to=2-2]
  \end{tikzcd}\]
  (2) $\implies$ (1) is clear.
\end{proof}

\section*{The split opfibration factorisation system}
% \label{sec:awfs}

There are algebraic formulations of normal and split opfibrations,
which we make use of in our type theoretic analysis.
These considerations lead us to
an algebraic weak factorisation system on $\Cat$,
where the algebras are precisely split opfibrations.

Any functor $F : \DD \to \CC$ factors through the
comma category as follows.
\[ \begin{tikzcd}
  \DD && \CC \\
  & {\FF / \CC}
  \arrow["F", from=1-1, to=1-3]
  \arrow["{L_F}"', from=1-1, to=2-2]
  \arrow["{R_F}"', from=2-2, to=1-3]
\end{tikzcd} \]
The factorisation of $F$ can be given by extending the pullback-hom
diagram (\cref{eq:pullback-hom})
\[
\begin{tikzcd}
  \DD & {\DD^\two} && \\
  && {F / \CC} & \DD \\
  \CC && {\CC^\two} & \CC
  \arrow["\Delta", from=1-1, to=1-2]
  % \arrow["{L_F}"{description, pos=0.3}, bend right, dashed, from=1-1, to=2-3]
  \arrow["F"', from=1-1, to=3-1]
  \arrow["{h_F}"{description}, from=1-2, to=2-3]
  \arrow["\src", bend left, from=1-2, to=2-4]
  \arrow["\pi_\DD", from=2-3, to=2-4]
  % \arrow["{R_F}"{description}, dashed, from=2-3, to=3-1]
  \arrow["{\pi_{\CC^\two}}"', from=2-3, to=3-3]
  \arrow["\lrcorner"{anchor=center, pos=0.125}, draw=none, from=2-3, to=3-4]
  \arrow["F", from=2-4, to=3-4]
  \arrow["\trg", from=3-3, to=3-1]
  \arrow["\src"', from=3-3, to=3-4]
  \arrow["F^\two"{description}, bend right, from=1-2, to=3-3]
\end{tikzcd}
\]
so that $L_F = h_F \circ \Delta : \DD \to F / \CC$ is as defined in
\Cref{prop:normal-opfibration-equiv-defs}
and $R_F = \trg \circ \pi_{\CC^\two} : F / \CC \to \CC$ is the codomain projection.
It follows that $R_F$ and $L_F$ factor $F$.
\[R_F \circ L_F = \trg \circ F^\two \circ \Delta = \trg \circ \Delta \circ F = F\]

\begin{proposition}\label{prop:normal-opfibration-algebra}
  Let $F : \DD \to \CC$ be a functor.
  The following are equivalent (structures):
  \begin{itemize}
    \item A normal opfibration $(F,l)$.
    \item A chosen left adjoint $\alpha \dashv L_F : \DD \to F / \CC$
    such that the counit of the adjunction is the identity
    \[\alpha \circ L_F = \id_\DD\]
    (making $\alpha$ ``left-adjoint-left-inverse'' or ``lali'')
    and $F \circ \alpha = R_F$, i.e. the adjunction is over $\CC$.
  \end{itemize}
\end{proposition}
\begin{proof}
  Assuming a normal opfibration $(F,l)$,
  consider the following diagram
  \[ \begin{tikzcd}
    \DD & {\DD^\two} & {F / \CC}
    \arrow[""{name=0, anchor=center, inner sep=0}, "\Delta", shift left=2, from=1-1, to=1-2]
    \arrow["{L_F}"{description}, shift left=3, bend left, from=1-1, to=1-3]
    \arrow[""{name=1, anchor=center, inner sep=0}, "\trg", shift left=2, from=1-2, to=1-1]
    \arrow[""{name=2, anchor=center, inner sep=0}, "{h_F}", shift left=2, from=1-2, to=1-3]
    \arrow["\alpha"{description}, shift left=3, bend left, from=1-3, to=1-1]
    \arrow[""{name=3, anchor=center, inner sep=0}, "l", shift left=2, from=1-3, to=1-2]
    \arrow["\dashv"{anchor=center, rotate=90}, draw=none, from=1, to=0]
    \arrow["\dashv"{anchor=center, rotate=90}, draw=none, from=3, to=2]
  \end{tikzcd}\]
  The composition of two left adjoints provides $\alpha = \trg \circ l$.
  The counit $\eta''_d : \alpha L_F d \to d$,
  for some $d \in \DD$,
  can be computed in terms of the respective counits $\eta$ and $\eta'$ of
  $\trg \dashv \Delta$ and $l \dashv h_F$.
  \[ \eta''_d = \eta_d \circ \trg (\eta'_{\Delta d}) : \alpha L_F d \to \trg \Delta d \to d \]
  It follows from normality of $l$ that
  $\trg (\eta'_{\Delta d})$ is the identity.
  On the other hand, $\eta_d$ is always the identity by a straightforward computation.
  Hence $\alpha$ is lali.
  Since $l$ is lari, 
  \[F \circ \alpha = F \circ \trg \circ l = \trg \circ F^\two \circ l = R_F \circ h_F \circ l = R_F\]

  To sketch the converse,
  suppose $\alpha$ is left-adjoint-left-inverse to $L_F$ such that
  $F \circ \alpha = R_F$.
  We want to construct a functor $l : F / \CC \to \DD^\two$ 
  fitting into the diagram above.
  % \[\begin{tikzcd}
  %   \DD & {\DD^\two} & {F / \CC}
  %   \arrow[""{name=0, anchor=center, inner sep=0}, "\Delta", shift left=2, from=1-1, to=1-2]
  %   \arrow["{L_F}"{description}, shift left=3, bend left, from=1-1, to=1-3]
  %   \arrow[""{name=1, anchor=center, inner sep=0}, "\trg", shift left=2, from=1-2, to=1-1]
  %   \arrow[""{name=2, anchor=center, inner sep=0}, "{h_F}", shift left=2, from=1-2, to=1-3]
  %   \arrow["\alpha"{description}, shift left=3, bend left, from=1-3, to=1-1]
  %   \arrow[""{name=3, anchor=center, inner sep=0}, "l", shift left=2, dashed, from=1-3, to=1-2]
  %   \arrow["\dashv"{anchor=center, rotate=90}, draw=none, from=1, to=0]
  %   \arrow["\dashv"{anchor=center, rotate=90}, draw=none, from=3, to=2]
  % \end{tikzcd}\]
  Note that $L_F$ also has a right adjoint,
  the projection $\pi_\DD : F / \CC \to \DD$ (in fact $L_F$ is lari).
  For an object $X = (x,y,f : Fx \to y)$ in $F / \CC$,
  let us write $\epsilon_X : L_F x \to X$
  for the counit of the adjunction $L_F \dashv \pi_\DD : F / \CC \to \DD$.
  To define $l : F / \CC \to \DD^\two$ on objects,
  % we apply $\alpha$ to this counit:
  take $X = (x,y,f : Fx \to y)$ in $F / \CC$ to 
  the object (an arrow in $\DD$) $\alpha \epsilon_X$ in $\DD^\two$.
\end{proof}
In fact, there is a monad with its underlying functor
$R : \Cat^\two \to \Cat^\two$
taking $F \mapsto R_F$.
The second condition of \Cref{prop:normal-opfibration-algebra}
implies that $F$ has a
pointed endofunctor algebra structure.
(The converse is not true;
a pointed endofunctor algebra provides normal lifts,
but those lifts may not be opcartesian.
In this sense, the pointed endofunctor algebras are equivalent to
normal Hurewicz maps \Cref{def:hurewicz-defs}.)
In the following diagram,
the unit is given by the left square,
and the algebra map is given by the right square.
\[
\begin{tikzcd}
  \DD & {F / \CC} & \DD \\
  \CC & \CC & \CC
  \arrow["{L_F}", from=1-1, to=1-2]
  \arrow["F"', from=1-1, to=2-1]
  \arrow["\alpha", from=1-2, to=1-3]
  \arrow["{R_F}"', from=1-2, to=2-2]
  \arrow["F", from=1-3, to=2-3]
  \arrow[equals, from=2-1, to=2-2]
  \arrow[equals, from=2-2, to=2-3]
\end{tikzcd}
\]
Next, we consider actual monad algebras and split opfibrations.

The following theorem can be found in \cite[2.13]{garner2008}. % and \cite{bourke2016}
A version of the theorem using pseudoalgebras rather than
algebras can be found in \cite{street1974},
where cloven opfibrations replace split opfibrations.
Recall the definition of an algebraic weak factorisation system (AWFS),
or ``natural weak factorisation system'' in \cite{garner2008}.
\begin{theorem}\label{thm:lari-opfib-awfs}
  The mapping $F \mapsto L_F$ defines a comonad
  $L : \Cat^\two \to \Cat^\two$
  and $F \mapsto R_F$ defines a monad $R : \Cat^\two \to \Cat^\two$.
  This forms the basis of
  an AWFS on $\Cat$.
\end{theorem}

\begin{definition}\label{def:morphism-split-opfibration}
A morphism of split opfibrations
(also known as an opcartesian functor) between
$(F : \DD \to \CC,l)$ and $(F' : \DD' \to \CC',l')$
consists of a functor $H : \CC \to \CC'$,
a functor $G : \DD \to \DD'$,
such that $F' \circ G = H \circ F$
and $G^\two \circ l = l' \circ G / H$.
\[ 
  \begin{tikzcd}
    \DD & {\DD'} \\
    \CC & \CC'
    \arrow["G", from=1-1, to=1-2]
    \arrow["F"', from=1-1, to=2-1]
    \arrow["H"', from=2-1, to=2-2]
    \arrow["{F'}", from=1-2, to=2-2]
  \end{tikzcd}
  \quad
  \quad \begin{tikzcd}
	{\DD^\two} & {\DD'^\two} \\
	{F / \CC} & {F' / \CC'}
	\arrow["{G^\two}", from=1-1, to=1-2]
	\arrow["l", from=2-1, to=1-1]
	\arrow["{G / H}"', from=2-1, to=2-2]
	\arrow["{l'}"', from=2-2, to=1-2]
\end{tikzcd} \]
Since $\FF$-opcartesian maps in $\DD$ are
isomorphic to the lifts of their images given by $l$,
the second condition implies that
$G$ takes all opcartesian maps over $\CC$ to opcartesian maps over $\CC'$.
\end{definition}

Let us write $\opfibcart$ for the category with
split opfibrations as objects,
and morphisms of split opfibrations as morphisms.

We also write $\lari$ for the category that has lari adjunctions $L \dashv R$
as objects -- meaning that the unit is the identity --
with morphisms of lari adjunctions defined as follows.
Suppose $L : A \to B$ and $L' : A' \to B'$ are lari,
with rights adjoints $R$ and $R'$ respectively,
a morphism $(F,G) : (L \dashv R) \to (L' \dashv R')$
consists of functors $F : A \to A'$ and $G : B \to B'$
such that $L' F = G L$, $R' G = F R$,
and for any object $b$ in $B$,
the two counits $\epsilon$ and $\epsilon'$ satisfy
$G \epsilon_b = \epsilon'_{G b}$.
\[
  \begin{tikzcd}
	B & {B'} \\
	A & {A'} \\
	B & {B'}
	\arrow["G", from=1-1, to=1-2]
	\arrow["R"', from=1-1, to=2-1]
	\arrow[""{name=0, anchor=center, inner sep=0}, "{=}"', shift right = 5, bend right = 60, from=1-1, to=3-1]
	\arrow["{R'}", from=1-2, to=2-2]
	\arrow[""{name=1, anchor=center, inner sep=0}, "{=}", shift left = 5, bend left = 60, from=1-2, to=3-2]
	\arrow["F"', from=2-1, to=2-2]
	\arrow["L"', from=2-1, to=3-1]
	\arrow["{L'}", from=2-2, to=3-2]
	\arrow["G"', from=3-1, to=3-2]
	\arrow["\epsilon", Rightarrow, from=2-1, to=0, shorten >=0.5ex]
	\arrow["{\epsilon'}"', Rightarrow, from=2-2, to=1, shorten >=0.5ex]
\end{tikzcd}
\]

\begin{theorem}
  For a functor $F : \DD \to \CC$,
  the category of algebras of the monad $R : \Cat^\two \to \Cat^\two$
  (from \Cref{thm:lari-opfib-awfs})
  is equivalent to the category of split opfibrations.
  The category of coalgebras of the comonad $L : \Cat^\two \to \Cat^\two$
  is equivalent to the category of lari adjunctions.
  \[ \alg{R} \simeq \opfibcart
  \quad \quad \quad \quad
  \coalg{L} \simeq \lari\]
  \[
  (F, \alpha : R_F \to F) \mapsto (F,l)
  \quad \quad
  (F, \alpha : F \to L_F) \mapsto F \dashv A
  \]

\end{theorem}

% \begin{theorem}
%   For a functor $F : \DD \to \CC$,
%   an algebra $(F, \alpha : R_F \to F)$ of the monad $R : \Cat^\two \to \Cat^\two$
%   of \Cref{thm:lari-opfib-awfs}
%   is equivalent (as structure) to a split opfibration $(F,l)$.
%   A coalgebra $(F, \alpha : F \to L_F)$ of the comonad $L : \Cat^\two \to \Cat^\two$
%   is equivalent to an adjunction $F \dashv A$
%   in which $F$ is lari.
%   \[ \alg{R} \simeq \opfibcart
%   \quad \quad
%   \coalg{L} \simeq \lari\]
% \end{theorem}
% \begin{proof}
  % We sketch the construction of a split opfibration from a monad algebra.
  % Given a monad algebra $(F : \DD \to \CC, {\alpha} : F / \CC \to \DD)$ of $R$,
  % we obtain a left adjoint $\alpha \dashv L_F$
  % of the left-map $L_F : \DD \to \CC$
  % that is also a pointed endofunctor algebra structure on $F$.
  % By \Cref{prop:normal-opfibration-algebra},
  % this provides a normal opfibration.
  % One can check that the lift $l$ defined in
  % the proposition is also split.
% \end{proof}


\section*{Straightening and unstraightening}

\Cref{thm:opfibration-main} is the standard straightening-unstraightening result
for opfibrations.
In preparation for the theorem,
we define a category consisting of
split opfibrations over a category $\CC$.
We write $\opfib(\CC)$ for the full subcategory of the slice
$\Cat / \CC$ consisting of split opfibrations.
We write $\opfibcart(\CC)$ for
the wide subcategory of $\opfib(\CC)$
consisting of morphisms of split opfibrations in the slice.
% One candidate for this is the full subcategory of the slice,
% i.e. the objects are split opfibrations,
% with no restrictions on the morphisms.
% It turns out that a better behaved category
% is given when we restrict to morphisms
% of split opfibrations in the slice,
That is, morphisms
\[(G,H) : (F : \DD \to \CC, l) \to (F' : \DD' \to \CC, l')\]
such that $H$ is the identity functor on $\CC$
(see \cite[Definition 3.1]{streicher2018}).
% We write $\opfibcart(\CC)$ for
% the wide subcategory of $\opfib(\CC)$
% consisting of morphisms of split opfibrations in the slice.
% \[ \smlopfibcart(\CC) \to \smlopfib(\CC)\]

% \begin{example}
%   We provide an example of a morphism in $\opfib(\Set)$
%   that is not in $\opfibcart(\Set)$.
%   That is, a morphism that does not preserve lifts.
%   Since $\Set$ has all pushouts,
%   the domain map $\src : \Set^\two \to \Set$ is an opfibration,
%   we can even make explicit choices of pushouts to make it split.

% \end{example}

\begin{theorem}\label{thm:opfibration-main}
  For any category $\CC$,
  the category of functors and natural transformations
  $[\CC,\Cat]$ is equivalent to $\opfibcart(\CC)$.
  \[[\CC,\Cat] \simeq \opfibcart(\CC)\] 
\end{theorem}
\begin{proof}
  % See for example \cite{streicher2018}.
  We introduce some notation,
  and sketch the functors that make up the equivalence.
   
  Recall the \emph{Grothendieck construction} on $A$,
  denoted $\int A$:
  objects of $\int A$ consist of pairs $(x, a)$,
  where $x \in \CC$ and $a \in A(x)$;
  morphisms $(x,a) \to (y,b)$ consist of pairs $(f,\phi)$,
  where $f : x \to y$ is a morphism in $\CC$
  and $\phi : A f a \to b$ is a morphism in $A(y)$.
  There is an associated forgetful functor $u_A : \int A \to \CC$,
  which takes the first component on objects and morphisms.
  The forgetful functor $u_A : \int A \to \CC$
  has a canonical split opfibration structure.
  Overall, this makes up functor $\int : [\CC,\Cat] \to \opfibcart(\CC)$.

  Conversely, for a split opfibration $(F : \DD \to \CC, l)$ over $\CC$,
  we can construct a functor $F^* : \CC \to \Cat$ by taking
  $x$ in $\CC$ to the fibre $F^* x$.
  Then the lifting operation $l$ provides, for each $f : x \to y$ in $\CC$,
  a functor $F^* f : F^* x \to F^* y$.
\end{proof}

A version of the following lemma can be found, for example,
in \cite[Theorem 4.1]{streicher2018}.
Recall that a discrete opfibration is an opfibration
with fibres that are sets.
\begin{lemma}\label{lem:split-opfib-between-split-opfibs}
  Let $(A \to X, \alpha)$ and $(B \to X, \beta)$ be two
  split opfibrations and $F : A \to B$
  a morphism of split opfibrations over $X$.
  Then $F$ is a discrete opfibration
  if and only if 
  over each $x$ in $X$,
  the functor between fibres $F^* x : A^* x \to B^* x$
  is a discrete opfibration.

  More generally,
  suppose that $F^* : X \to \Cat^\two$ factors through the
  category of split fibrations in $\Cat$.
  \[
    \begin{tikzcd}
      & \opfibcart \\
      X & {\Cat^\two}
      \arrow[from=1-2, to=2-2]
      \arrow["{(F,\phi)}", dashed, from=2-1, to=1-2]
      \arrow["F"', from=2-1, to=2-2]
    \end{tikzcd}
  \]
  This means that for each $x$ in $X$,
  the functor between fibres $F^* x : A^* x \to B^* x$
  is equipped with the structure of a split opfibration
  $(F^* x, \phi_x)$,
  and for each $t : x \to y$ in $X$,
  transporting along $t$ forms
  a morphism of split opfibrations:
  \[(A^* t, B^* t) : (F^* x,\phi_x) \to (F^* y,\phi_y)
    \quad \quad \quad
    \begin{tikzcd}
      {A^* x} & {A^* y} \\
      {B^* x} & {B^* y}
      \arrow["{A^* t}", from=1-1, to=1-2]
      \arrow["{(F^* x,\phi_x)}"', from=1-1, to=2-1]
      \arrow["{(F^* y,\phi_y)}", from=1-2, to=2-2]
      \arrow["{B^* t}"', from=2-1, to=2-2]
    \end{tikzcd} \]
  Then we have a split opfibration structure on $F$.
\end{lemma}
\begin{proof}
  We define a lifting operation $\psi$ for $F$,
  suppose $a \in A$ and $f : F a \to b$.
  Let us write $t : x \to y$ for the image of $f : F a \to b$ in $X$.
  We can lift $t : x \to y$ using $\alpha$ and $\beta$ to get 
  $\alpha_{(a,t)} : a \to a_y$ and $\beta_{(F a, t)} : F a \to F a_y$;
  we have $F \alpha_{(a,t)} = \beta_{(F a, t)}$,
  since $F$ is a morphism of split opfibrations.
  Since $\beta_{(F a, t)}$ is opcartesian, 
  there is a unique ``vertical'' map $u : F a \to b$ --
  meaning that it is in the fibre over $y$ --
  such that $f = u \circ \beta_{(F a, t)}$.
  We lift $u$, using $\phi_y$ to obtain a morphism
  $\phi_{(a_y,u)} : a_y \to a_y'$ in the fibre of $A$ over $y$.
  \[\begin{tikzcd}
    a & {a_y} & {a_y'} \\
    \\
    Fa & {F a_y} & b \\
    \\
    x & y & y
    \arrow["{\alpha_{(a,t)}}", from=1-1, to=1-2]
    \arrow["{\psi_{(a,f)}}"', bend right, from=1-1, to=1-3]
    \arrow["{\phi_{(a_y,u)}}", from=1-2, to=1-3]
    \arrow["{\beta_{(F a, t)}}", from=3-1, to=3-2]
    \arrow["f"', bend right, from=3-1, to=3-3]
    \arrow["u", from=3-2, to=3-3]
    \arrow["t", from=5-1, to=5-2]
    \arrow[equals, from=5-2, to=5-3]
  \end{tikzcd}\]
  We define $\psi_{(a,f)}$ to be the composition
  \[\psi_{(a,f)} = u \circ \alpha_{(a,t)} : a \to a_y'\]

  It follows from this definition
  and normality of $\alpha$, $\beta$ and $\phi_y$
  that if $f$ is the identity,
  then so are $t$, $\beta_{(F a, t)}$, $u$, $\alpha_{(a,t)}$,
  and ${\phi_{(a_y,u)}}$.
  Hence $\psi$ is normal.

  To check that $\psi$ is split,
  suppose additionally $g : b \to b'$.
  Let us denote the image of $g$ in $X$ as $s : y \to z$.
  We can construct a lift $\psi_{(a'_y,g)} : a'_y \to a''_z$
  of $g$ along $F$,
  just as we constructed the lift of $f$,
  as the composition of the lift $\alpha_{(a'_y,s)} : a'_y \to a'_z$
  of $s$ followed by the lift $\phi_{(a_z,v)} : a'_z \to a''_z$ of 
  $v : F a'_z \to F a''_z$, 
  the unique map over $z$ satisfying
  $v \circ \beta_{(b,s)} = g$
  (refer to the diagram below, where the situation is similar).
  We write $\alpha_{(a_y,s)} : a_y \to a_z$,
  which is sent by $F$ to $\beta_{(F a_y,s)} : F a_y \to F a_z$ and
  $u' : F a_z \to F a'_z$ for the unique map over $z$
  satisfying $u' \circ \beta_{(F a_y,s)} = \beta_{(b,s)} \circ u$.
  The condition that
  $(A^* s, B^* s) : (F^* y,\phi_y) \to (F^* z,\phi_z)$
  is a morphism of split opfibrations can be described in
  more elementary terms as a coherence condition:
  \[\phi_{(a_z,u')} \circ \alpha_{(a_y,s)} = \alpha_{(a'_y,s)} \circ \phi_{(a_y,u)}\]
  \[\begin{tikzcd}
    &&& {a_y} & {a_z} \\
    {A^* y} & {A^* z} && {a_y'} & {a'_z} \\
    {B^* y} & {B^* z} && {F a_y} & {F a_z} \\
    y & z && b & {F a'_z}
    \arrow["{{{{\alpha_{(a_y,s)}}}}}", from=1-4, to=1-5]
    \arrow["{{{{\phi_{(a_y,u)}}}}}"', from=1-4, to=2-4]
    \arrow["{{{{\phi_{(a_z,u')}}}}}", from=1-5, to=2-5]
    \arrow["{{A^*s}}", from=2-1, to=2-2]
    \arrow["{{(F^* y,\phi_y)}}"', from=2-1, to=3-1]
    \arrow["{{(F^* z,\phi_z)}}", from=2-2, to=3-2]
    \arrow["{{{{\alpha_{(a'_y,s)}}}}}"', from=2-4, to=2-5]
    \arrow["{{B^*s}}"', from=3-1, to=3-2]
    \arrow["{{{{\beta_{(F a_y, s)}}}}}", from=3-4, to=3-5]
    \arrow["u"', from=3-4, to=4-4]
    \arrow["{{{u'}}}", from=3-5, to=4-5]
    \arrow["s", from=4-1, to=4-2]
    \arrow["{{{{\beta_{(b, s)}}}}}"', from=4-4, to=4-5]
  \end{tikzcd}\]
  Since $\phi_z$, $\alpha$ and $\beta$ are split,
  \begin{align*}
    & \psi_{(a, gf)} \\
  = \, & \phi_{(a_z,v u')} \circ \alpha_{(a,st)}\\
  = \, & \phi_{(a'_z,v)} \circ \phi_{(a_z,u')} \circ \alpha_{(a_y,s)} \circ \alpha_{(a,t)} \\
  = \, & \phi_{(a'_z,v)} \circ \alpha_{(a'_y,s)} \circ \phi_{(a_y,u)} \circ \alpha_{(a,t)} \\
  = \, & \psi_{(a'_y,g)} \circ \psi_{(a,f)}
  \end{align*}
  
  We show that $\psi_{(a,f)}$ is opcartesian over $B$.
  Suppose $h : a \to e$ is a morphism in $A$
  and $g : b \to F e$ such that $F h = g \circ f$.
  We want to show that there is a unique lift
  $g' : a'_y \to e$ of $g$ in $A$
  such that $g' \circ \psi_{(a,f)} = h$.
  Let us reuse the notation for maps $u$, $s$, $v$ and $u'$ from above,
  as well as the names for the endpoints of lifts.
  Since $\alpha_{(a,st)} : a \to a_z$ is opcartesian over $X$,
  there is a unique map $w : a_z \to e$
  over $z$ satisfying $w \circ \alpha_{(a,st)} = h$.  
  Since $\phi_{(a_z,vu')} : a_z \to a''_z$ is opcartesian over $B^* z$,
  there is a unique map $w' : a''_z \to e$
  over $F e$ satisfying $w' \circ \phi_{(a_z,vu')} = w$.
  Clearly, $F (w' \circ \psi_{(a'_y,g)}) = g$ gives us a lift.
  \[\begin{tikzcd}[column sep = large]
	&& {a_z} &&&& {F a_z} \\
	& {a_y} & {a'_z} &&& {F a_y} & {F a'_z} \\
	a & {a_y'} & {a''_z} && Fa & b & {F a''_z} \\
	&& e &&&& {F e} \\
	x & y & z && x & y & z
	\arrow["{{\phi_{(a_z,u')}}}", from=1-3, to=2-3]
	\arrow["{u'}", dashed, from=1-7, to=2-7]
	\arrow["{{\alpha_{(a_y,s)}}}", from=2-2, to=1-3]
	\arrow[from=2-2, to=2-3]
	\arrow["{{\phi_{(a_y,u)}}}"{description}, from=2-2, to=3-2]
	\arrow["{{\phi_{(a'_z,v)}}}", from=2-3, to=3-3]
	\arrow["{{\beta_{(F a_y, s)}}}", from=2-6, to=1-7]
	\arrow[from=2-6, to=2-7]
	\arrow["u"{description}, dashed, from=2-6, to=3-6]
	\arrow["v", dashed, from=2-7, to=3-7]
	\arrow["{{\alpha_{(a,t)}}}", from=3-1, to=2-2]
	\arrow["{\psi_{(a,f)}}"', from=3-1, to=3-2]
	\arrow["h"', bend right, from=3-1, to=4-3]
	\arrow["{{\alpha_{(a'_y,s)}}}"{description}, from=3-2, to=2-3]
	\arrow["{\psi_{(a'_y,g)}}"', from=3-2, to=3-3]
	\arrow["{g'}"', bend right = 10, dashed, from=3-2, to=4-3]
	\arrow["w'", dashed, from=3-3, to=4-3]
	\arrow["{{\beta_{(F a, t)}}}", from=3-5, to=2-6]
	\arrow["f"{description}, from=3-5, to=3-6]
	\arrow["{F h}"', bend right, from=3-5, to=4-7]
	\arrow["{{\beta_{(b, s)}}}"{description}, from=3-6, to=2-7]
	\arrow["g"', from=3-6, to=3-7]
	\arrow["g"', bend right = 10, from=3-6, to=4-7]
	\arrow[equals, from=3-7, to=4-7]
	\arrow["t", from=5-1, to=5-2]
	\arrow["s", from=5-2, to=5-3]
	\arrow["t", from=5-5, to=5-6]
	\arrow["s", from=5-6, to=5-7]
  \arrow["w"{description}, shift left=3, bend left = 70, dashed, from=1-3, to=4-3]
  \end{tikzcd}\]
  Let $g' : a'_y \to e$ be such that $F g' = g$.
  It suffices to show that $g' \circ \psi_{(a,f)} = h$
  if and only if $g' = w' \circ \psi_{(a'_y,g)}$.
  It is easy to check that $\phi_{(a_y, u)}$ is not just
  opcartesian with respect to $F^* y : A^* x \to B^* x$,
  but also opcartesian with respect to all of $F : A \to B$.
  Also, $\alpha_{(a,t)}$ is opcartesian over $X$.
  Thus, 
  \begin{align*}
    & g' = w' \circ \psi_{(a'_y,g)} \\
    \iff \, & g' \circ \phi_{(a_y, u)} = w' \circ \psi_{(a'_y,g)} \circ \phi_{(a_y, u)} \\
    \iff \, & g' \circ \phi_{(a_y, u)} \circ \alpha_{(a,t)} = w' \circ \psi_{(a'_y,g)} \circ \phi_{(a_y, u)} \circ \alpha_{(a,t)} \\
    \iff \, & g' \circ \psi_{(a,f)} = w' \circ \psi_{(a'_y,g)} \circ \psi_{(a,f)} \\
  \end{align*}
  Finally,
  \[ w' \circ \psi_{(a'_y,g)} \circ \psi_{(a,f)} = w' \circ \psi_{(a,gf)}
  = w' \circ \phi_{(a_z,vu')} \circ \alpha_{(a,st)} = w \circ \alpha_{(a,st)} = h \]

  Now suppose instead that each $F^* x : A^* x \to B^* x$ is a discrete opfibration.
  This is the same as saying $F^* x : A^* x \to B^* x$ is a split opfibration
  and for each $b \in B^* x$ the fibre
  $(F^* x)^* b \iso F^* b$ is a set.
  Let us first assume the coherence condition for transports holds.
  Then we have that $F$ is a split opfibration such that for any $b \in B$,
  $F^* b$ is a set.
  In other words, $F$ is a discrete opfibration.

  Finally, we verify the coherence condition for transports.
  Suppose $s : y \to z$ in $X$, $a_y \in A$ is over $y$, and $u : F a_y \to b$ in $B$.
  We generate a similar diagram as before:
  \[
  \begin{tikzcd}
    {a_y} & {a_z} & \\
    {a_y'} & {a'_z} & {a''_z} \\
    {F a_y} & {F a_z} \\
    b & {F a'_z} \\
    y & z
    \arrow["{{{{\alpha_{(a_y,s)}}}}}", from=1-1, to=1-2]
    \arrow["{{{{\phi_{(a_y,u)}}}}}"', from=1-1, to=2-1]
    \arrow["v", dashed, from=1-2, to=2-2]
    \arrow["{{{{\phi_{(a_z,u')}}}}}", from=1-2, to=2-3]
    \arrow["{{{{\alpha_{(a'_y,s)}}}}}"', from=2-1, to=2-2]
    \arrow["{{{{\beta_{(F a_y, s)}}}}}", from=3-1, to=3-2]
    \arrow["u"', from=3-1, to=4-1]
    \arrow["{{{u'}}}", from=3-2, to=4-2]
    \arrow["{{{{\beta_{(b, s)}}}}}"', from=4-1, to=4-2]
    \arrow["s", from=5-1, to=5-2]
  \end{tikzcd}
  \]
  where $u' : F a_z \to F a'_z$ is the unique map over $z$ such that
  \[u' \circ \beta_{(F a_y, s)} = \beta_{(b,s)} \circ u\]
  and $v : a_z \to a'_z$ is the unique map over $z$ such that
  \[v \circ \alpha_{(a_y,s)} = \alpha_{(a'_y,s)} \circ \phi_{(a_y,u)}\]
  It suffices to show that $v = \phi_{(a_z,u')}$.
  Since $F^* z$ is a discrete opfibration
  and $v$ and $\phi_{(a_z,u')}$ are both in the fibre over $z$,
  this amounts to showing that $F v = u'$.
  By definition of $u'$, we simply need to check the following.
  \[
    F v \circ \beta_{(F a_y, s)}
    = F (v \circ \alpha_{(a_y,s)})
    = F (\alpha_{(a'_y,s)} \circ \phi_{(a_y,u)})
    = \beta_{(b,s)} \circ u
  \]
\end{proof}

There is a similar result for lari functors between split opfibrations.
\begin{lemma} \label{lem:lari-between-split-opfibs}
  Let $(A \to X, \alpha)$ and $(B \to X, \beta)$ be two
  split opfibrations and $F : A \to B$
  a morphism of split opfibrations over $X$.
  Suppose that $F^* : X \to \Cat^\two$ factors through the
  category of lari adjunctions in $\Cat$.
  \[
    \begin{tikzcd}
      & \lari \\
      X & {\Cat^\two}
      \arrow[from=1-2, to=2-2]
      \arrow["{(F,R)}", dashed, from=2-1, to=1-2]
      \arrow["F^*"', from=2-1, to=2-2]
    \end{tikzcd}
  \]
  Then we have a lari adjunction $F \dashv R$ such that
  the right adjoint $R$ is a morphism of split opfibrations over $X$.
\end{lemma}
\begin{proof}
  We first unfold the assumption that $F^*$ factors through $\lari$.
  This means that for each $x$ in $X$,
  the functor between fibres $F^* x : A^* x \to B^* x$
  is equipped with a right adjoint $R_x : B^* x \to A^* x$,
  such that the unit is the identity.
  Also, for each $t : x \to y$ in $X$,
  % transporting along $t$ is a morphism of lari functors
  % $(A^* t, B^* t) : (F^* x \dashv R_x) \to (F^* y \dashv R_y)$.
  % Unfolding this further means that
  $F^* y \circ A^* t = B^* t \circ F^* x$, $R_y \circ B^* t = A^* t \circ R_x$,
  and for any object $b$ in $B^* x$,
  the counits satisfy
  $B^* t \epsilon_b = \epsilon'_{B^* t b}$.
  
  To define the functor $R : B \to A$,
  let $R$ take an object $b \in B$ over $x \in X$ to $R_x b$.
  For a morphism $f: b \to b'$ over $t : x \to y$,
  we take the canonical opcartesian-vertical factorisation
  \[ f = u \circ \beta_{(b,t)}\]
  for $\beta_{(b,t)} : b \to b_y$
  and $u : b_y \to b'$ over $y$.
  We take $R f := R_y u \circ \alpha_{(R_x b,t)}$,
  checking that the domain of $R_y u$
  is equal to the codomain of $\alpha_{(R_x b,t)}$,
  \[R_y b_y = R_y (B^* t (b)) = A^* t (R_x b)\]
  \[\begin{tikzcd}
    & {R_x b} && b & x \\
    {R_y b_y} & {R_y b'} & {b_y} & {b'} & y
    \arrow["{\alpha_{(R_x b,t)}}"', from=1-2, to=2-1]
    \arrow["{R f}", dashed, from=1-2, to=2-2]
    \arrow["{\beta_{(b,t)}}"', from=1-4, to=2-3]
    \arrow["f", from=1-4, to=2-4]
    \arrow["t"', from=1-5, to=2-5]
    \arrow["{R_y u}"', from=2-1, to=2-2]
    \arrow["u"', from=2-3, to=2-4]
  \end{tikzcd}\]
  It is easy to check that $R$ is functorial.
  Furthermore, it is evidently a morphism of split opfibrations over $X$,
  so that for $b \in B$ over $x \in X$ and $t : x \to y$,
  \[R \beta_{(b,t)} = \alpha{(R b,t)}\]
  and for $u : b \to b'$ vertical over $x$,
  \[R u = R_x u\]

  To make a lari adjunction $L \dashv R$,
  we check that $\id_A = R \circ F$ and construct a counit
  $\epsilon : F \circ R \to \id_B$ satisfying two equations
  $\epsilon F = \id$ and $R \epsilon = \id$.
  Firstly, for an object $a \in A$ over $x \in X$,
  \[R F a = R_x F_x a = a\]
  For morphisms, by taking opcartesian-vertical factorisations,
  it suffices to check that
  \[R F \alpha_{(a,t)} = R \beta_{(F a, t)} = \alpha_{(R F a,t)} = \alpha_{(a,t)}\]
  and for $u$ vertical over $y \in X$,
  \[R F u = R_y F_y u = u\]
  For an object $b \in B$ over $x \in X$,
  we define $\epsilon_b : F R b \to b$
  as the component of the counit of $F^* x \dashv R_x$ at $b$.
  To prove naturality of $\eta$ at a morphism
  $f : b \to b'$ over $t : x \to y$,
  factor $f = u \circ \beta_{(b,t)}$ as usual,
  where $\beta_{(b,t)} : b \to b_y$.
  Using the coherence equation, we have 
  \[B^* t \epsilon_b = \epsilon_{B^* t b} = \epsilon_{b_y}\]
  so that the following diagram commutes
  \[ \begin{tikzcd}
    {F_x R_x b} & b \\
    {F_y R_y b} & {b_y} \\
    {F_y R_y b} & {b'}
    \arrow["{\epsilon_{b}}", from=1-1, to=1-2]
    \arrow["{\beta_{(F R b,t)}}"', from=1-1, to=2-1]
    \arrow["{F R f}"', shift right=10, bend right = 40, from=1-1, to=3-1]
    \arrow["{\beta_{(b,t)}}", from=1-2, to=2-2]
    \arrow["f", shift left=5, bend left = 40, from=1-2, to=3-2]
    \arrow["{\epsilon_{b_y}}"', from=2-1, to=2-2]
    \arrow["{B^* t (\epsilon_b)}", draw=none, from=2-1, to=2-2]
    \arrow["{F_y R_y u}"', from=2-1, to=3-1]
    \arrow["u", from=2-2, to=3-2]
    \arrow["{\epsilon_{b'}}"', from=3-1, to=3-2]
  \end{tikzcd} \]
  Since each $\eta_b$ is defined as the counit of
  the adjunction $F^* x \dashv R_x$ over some $x$,
  it is vertical.
  The two remaining equations for $a \in A$ and $b \in B$ 
  $\epsilon_{F a} = \id_a$ and $R \epsilon_b = \id_{R b}$
  then follow automatically
  from each $F^* x \dashv R_x$ being lari.
\end{proof}

\section*{The universal split opfibration}
% \label{sec:universes-in-cat}

In this section, we describe the ``universal split opfibration'',
classifying \emph{small} split opfibrations in $\Cat$,
i.e. split opfibrations with small categories as fibres.
This appears, for example, in \cite{hess2007}
as the ``tautological bundle of categories''.

% We denote by $\cat$ the category of small categories,
% which is an object in $\Cat$.
We use $\smlopfibcart(\CC)$ to denote the category
of small split opfibrations over $\CC$
and morphisms of split opfibrations over $\CC$.
Let $\CC$ be a large category and
$A : \CC \to \cat$ be a functor.
Then \Cref{thm:opfibration-main}
restricts to the following equivalence.
\[[\CC,\cat] \simeq \smlopfibcart(\CC)\] 

Hess calls the following definition the
``tautological bundle of categories'' \cite{hess2007}.
Elementarily,
it is the forgetful functor from the 1-category of small pointed categories
to the 1-category of small categories.
We formulate the definition as a Grothendieck construction for convenience.

\begin{definition}\label{def:universal-small-split-opfibration}
  Consider the identity functor $\id_\cat : \cat \to \cat$.
  We write \[\pcat := \smallint \id_\cat\] for the Grothendieck construction
  on the identity functor
  and $\ulow := u_{\id_\cat} : \pcat \to \cat$ for the forgetful functor to the base.
  We call 
  % $\pcat$ the
  % \emph{(opfibrant) category of small pointed categories},
  % and
  $\ulow$ the \emph{universal small split opfibration}
  (or just the \emph{universal split opfibration}).
  This name will be justified by \Cref{prop:universal-opfibration-classifies}.
  We will also refer to $\ulow : \pcat \to \cat$ as a \emph{universe},
  as it will act like one in our model.
\end{definition}

The object $\pcat$ itself can be computed as follows:
it is the category where objects are pairs $(\CC,x)$ such that
$\CC$ is a small category and $x$ is an object in $\CC$,
and morphisms are pairs $(\FF,\phi) : (\CC,x) \to (\DD,y)$
such that $\FF : \CC \to \DD$ is a functor
and $\phi : \FF x \to y$ is a morphism in $\DD$.

For categories $\DD$ and $\CC$ and functors
$F : \DD \to \CC$ and $A : \CC \to \cat$,
there is a functor
$\int F : \int A \circ F \to \int A$
such that the following square is a pullback in $\Cat$.
\[
\begin{tikzcd}
  {\int A F} & {\int A} \\
  \DD & \CC 
  \arrow["{\int F}", from=1-1, to=1-2]
  \arrow["{u_{A F}}"', from=1-1, to=2-1]
  \arrow["{u_A}", from=1-2, to=2-2]
  \arrow["F"', from=2-1, to=2-2]
\end{tikzcd}
\]
It follows that the Grothendieck construction
is always a pullback of the universal one:
for any category $\CC$ and functor $A : \CC \to \cat$,
there is a pullback square in $\Cat$
\[
  \begin{tikzcd}
    {\int A} & {\pcat} \\
    \CC & \cat
    \arrow[from=1-1, to=1-2]
    \arrow["u_A"', from=1-1, to=2-1]
    \arrow["u_\cat", from=1-2, to=2-2]
    \arrow["A"', from=2-1, to=2-2]
  \end{tikzcd}
\]

Combining this with the equivalence
$[\CC,\cat] \simeq \smlopfibcart(\CC)$,
we have the following proposition.
\begin{proposition}\label{prop:universal-opfibration-classifies}
  Let $\CC$ and $\DD$ be large categories.
  Suppose $(F : \DD \to \CC, l)$ is a small split opfibration.
  Then $F$ is a pullback of the universal split opfibration
  $\ulow : \pcat \to \cat$ along its fibres functor
  $F^* : \CC \to \cat$.
\end{proposition}

The universal split opfibration
$\ulow : \pcat \to \cat$ is unique up to isomorphism,
but it is useful to distinguish its isomorphic ``vertical opposite''
which will be useful for our type theory.
Namely, we can consider the Grothendieck construction
on the opposite functor $\op : \cat \to \cat$.
\[ \catp := \smallint \op 
  \quad \quad \uupp := u_{\op{}} : \catp \to \cat \]
One can think of $\catp$ as $\pcat$,
except that we invert the
morphisms in the fibres.
% Since $\op{}$ is self-inverse, 
% these two universes are isomorphic.
\[\begin{tikzcd}
    \catp & \pcat \\
    \cat & \cat
    \arrow["\sim", from=1-1, to=1-2]
    \arrow["{\uupp}"', from=1-1, to=2-1]
    \arrow["\lrcorner"{anchor=center, pos=0.125}, draw=none, from=1-1, to=2-2]
    \arrow["{\ulow}", from=1-2, to=2-2]
    \arrow["\op{}"', from=2-1, to=2-2]
    \arrow["\sim", draw=none, from=2-1, to=2-2]
  \end{tikzcd} \]
More generally, viewing a split opfibration over $\CC$
as an \emph{indexed category} $A : \CC \to \cat$,
we can construct its ``vertical opposite'' by composing with the opposite functor
and taking the Grothendieck construction
$\int \op{} \circ A$,
which is to take the pullback of $\uupp$ along $A$.

Recall that a functor is a split fibration precisely when
its opposite is a split opfibration.
It follows that both
\[\vlow := {\uupp}^\op : {\catp}^\op \to {\cat}^\op
\quad \text{and} \quad
\vupp := {\ulow}^\op : {\pcat}^\op \to {\cat}^\op
\]
classify small split fibrations (they are isomorphic).
The universes $\vlow$ and $\ulow$ have opposite morphisms in their bases,
but parallel morphisms in their fibres,
because $\op$ reverses arrows in both the base and fibre.
Similarly $\vupp$ and $\uupp$ have opposite morphisms in their bases,
but parallel morphisms in their fibres.
We can also view $\vlow$ and $\vupp$ as classifying
\emph{contravariant indexed categories},
or \emph{$\cat$-valued presheaves} $A : \CC \to \cat^\op$.

We will call $\vlow : {\catp}^\op \to {\cat}^\op$
the \emph{universal split fibration},
since if we take the fiber of a category $\CC : 1 \to {\cat}^\op$
in $\vlow : {\pcat}^\op \to {\cat}^\op$,
we obtain $\CC$ itself,
whereas the fiber in $\vupp : {\pcat}^\op \to {\cat}^\op$ is $\CC^\op$.

\section*{Clans and exponentiability}

Let us write $\opfib$ for the class of split opfibrations
in $\Cat$ and $\smlopfib$ for the class of small split opfibrations.
We recall some standard facts about the two classes,
in preparation for type theoretic axiomatisation.

Both $\opfib$ and $\smlopfib$ are stable under pullback.
That $\smlopfib$ is stable under pullback follows from
\Cref{prop:universal-opfibration-classifies}.
One can check that both
$\opfib$ and $\smlopfib$ contain all isomorphisms
and are closed under composition.
Furthermore, for any category $\CC$,
the unique map $\CC \to 1$ to the terminal category
is a split opfibration,
which is small precisely when $\CC$ is.
Recalling the definition of a (pre)clan from \Cref{def:preclan},
we summarise these observations as follows.

\begin{proposition}\label{prop:preclans-in-cat}
  The category of large categories with the class of split opfibrations
  $(\Cat, \opfib)$ forms a clan.
  $\Cat$ with the class of small split opfibrations
  $(\Cat, \smlopfib)$ forms a preclan.
  
  Dually,
  $\Cat$ with the class of split fibrations
  $(\Cat, \fib)$ forms a clan
  and $\Cat$ with the class of small split fibrations
  $(\Cat, \smlfib)$ forms a preclan.
\end{proposition}

We turn our attention to the issue of exponentiability.
The failure of $\opfib$ and $\fib$ to be $\pi$-clans
distinguishes the type theory of categories from that of groupoids.
It is well known that opfibrations and fibrations are
exponentiable with respect to all maps (Conduch\'e)
\cite{giraud1964, conduche1972}.
Furthermore, pushforwards in $\Cat$ satisfy several closure conditions;
for the following proposition,
recall the definition of a $\tau$-preclan (\Cref{def:double-preclan}).

\begin{proposition}\label{prop:pushforward-in-cat}
  The opfibrations and fibrations in $\Cat$ form a $\tau$-clan
  $(\Cat,\opfib,\fib)$,
  and small opfibrations and small fibrations form a $\tau$-preclan
  $(\Cat,\smlopfib,\smlfib)$.
\end{proposition}
\begin{proof}
  By \Cref{lem:stable-class-of-exponentiable-maps},
  to show that $\fib \subseteq \Exp_\opfib$,
  it suffices to check that for any split opfibration $G : \EE \to \DD$
  and any split fibration $F : \DD \to \CC$,
  there is a split opfibration $P : \PP \to \CC$
  satisfying the universal property of the pushforward
  so that $F_* G = P$.
  Indeed, a pushforward can be constructed for
  any Conduch\'e functor \cite[Proposition 1.5.1]{vidmar2018},
  and fibrations are Conduch\'e \cite[Lemma 1.2.10]{vidmar2018}.
  In this case,
  since $G$ is a split opfibration and $F$ is a split fibration,
  one can extend the cited construction to produce a split opfibration
  $(F_* G, l)$.
\end{proof}

\begin{remark}\label{rmk:pi-preserve-opfib-mor}
  Given a fibration $f : Y \to X$,
  \Cref{prop:preclans-in-cat} and
  \Cref{prop:pushforward-in-cat} imply that there are functors
  between the full subcategories of the slice $\opfib(\star) \subseteq \Cat / (\star)$.
  \[
  \begin{tikzcd}[column sep = huge]
    {\opfib(Y)} & {\opfib(X)}
    \arrow[""{name=0, anchor=center, inner sep=0}, "{f_*}"', shift right=5, bend right, from=1-1, to=1-2]
    \arrow[""{name=1, anchor=center, inner sep=0}, "{f_!}", shift left=5, bend left, from=1-1, to=1-2]
    \arrow[""{name=2, anchor=center, inner sep=0}, "{f^*}"{description}, from=1-2, to=1-1]
    \arrow["\dashv"{anchor=center, rotate=-90}, draw=none, from=1, to=2]
    \arrow["\dashv"{anchor=center, rotate=-90}, draw=none, from=2, to=0]
  \end{tikzcd}
  \]
  In fact, it is also true that these functors restrict to the
  (non-full) subcategories of opfibrations and morphisms of opfibrations.
  \[
  \begin{tikzcd}
    {\opfibcart(Y)} & {\opfibcart(X)} \\
    {\opfib(Y)} & {\opfib(X)}
    \arrow["{f_!}", dashed, from=1-1, to=1-2]
    \arrow[from=1-1, to=2-1]
    \arrow[from=1-2, to=2-2]
    \arrow["{{f_!}}"', from=2-1, to=2-2]
  \end{tikzcd}
  \quad \quad 
  \begin{tikzcd}
    {\opfibcart(Y)} & {\opfibcart(X)} \\
    {\opfib(Y)} & {\opfib(X)}
    \arrow["{f_*}", dashed, from=1-1, to=1-2]
    \arrow[from=1-1, to=2-1]
    \arrow[from=1-2, to=2-2]
    \arrow["{{f_*}}"', from=2-1, to=2-2]
  \end{tikzcd}
  \quad \quad 
  \begin{tikzcd}
    {\opfibcart(Y)} & {\opfibcart(X)} \\
    {\opfib(Y)} & {\opfib(X)}
    \arrow["{f^*}"', dashed, from=1-2, to=1-1]
    \arrow[from=1-1, to=2-1]
    \arrow[from=1-2, to=2-2]
    \arrow["{{f^*}}", from=2-2, to=2-1]
  \end{tikzcd}
  \] 
  \[
  \begin{tikzcd}[column sep = huge]
    {\opfibcart(Y)} & {\opfibcart(X)}
    \arrow[""{name=0, anchor=center, inner sep=0}, "{f_*}"', shift right=5, bend right, from=1-1, to=1-2]
    \arrow[""{name=1, anchor=center, inner sep=0}, "{f_!}", shift left=5, bend left, from=1-1, to=1-2]
    \arrow[""{name=2, anchor=center, inner sep=0}, "{f^*}"{description}, from=1-2, to=1-1]
    \arrow["\dashv"{anchor=center, rotate=-90}, draw=none, from=1, to=2]
    \arrow["\dashv"{anchor=center, rotate=-90}, draw=none, from=2, to=0]
\end{tikzcd}\]
Unfortunately, the theory of exponentiability presented in \Cref{sec:exponentiable}
is inadequate for capturing this aspect of the model in $\Cat$.
One could naively throw the diagrams drawn above into the axiomatisation,
but this will not be enough,
as there should also be some extra coherence equations between
the 2-cells from the various related adjunctions.
A similar case is worked out in the language of comprehension categories in \cite{jacobs1993},
which could possibly be extended with the ideas from \Cref{sec:exponentiable}.
Related to this,
Gambino and Larrea construct models of Martin-L\"of Type Theory
from ``type theoretic'' AWFSs \cite{gambino2023},
which one could potentially modify to suit the case here.
Najmaei et al. present a type theory that internalises
the morphisms of algebras in an AWFS (as is the case here)
\cite{najmaei2026}.
\end{remark}

\section*{Twisted path factorisations}
% \label{sec:twisted-path-factorisation}
Motivated from the groupoid model \cite{hofmann1995},
and the algebraic semantics of identity types \cite{awodey2025}
and path types \cite{awodey2026}, 
we seek to produce a similar axiomatisation of
\emph{hom-types} in $\Cat$,
using Paige North's idea of two elimination rules \cite{north2019}
instead of the single one for identity types.
Recall from \cite{awodey2026} that in the category of groupoids $\Grpd$,
a (closed) type $A$ is interpreted as an isofibration (over $1$),
and the identity type for elements of $A$
is modelled by exponentiating $A$ by the bipointed
``walking isomorphism''.
\[ \begin{tikzcd}
	1 & I & A & {A^I} \\
	& {1+1} && {A \times A}
	\arrow["{!}"', from=1-2, to=1-1]
	\arrow[from=1-3, to=1-4]
	\arrow[from=1-3, to=2-4]
	\arrow["{{(\src, \trg)}}", from=1-4, to=2-4]
	\arrow["{!}", from=2-2, to=1-1]
	\arrow["{{(0,1)}}"', from=2-2, to=1-2]
\end{tikzcd} \]
For this to model a type,
we must ensure that the map $A^I \to A \times A$ is an isofibration,
which it is.

In $\Cat$, consider an opfibration over $1$, i.e. a category $A$.
It is natural to exponentiate by the
bipointed ``walking arrow'' $\two$, instead of $I$.
\[
\begin{tikzcd}
	1 & \two & A & {A^\two} \\
	& {1+1} && {A \times A}
	\arrow["{!}"', from=1-2, to=1-1]
	\arrow[from=1-3, to=1-4]
	\arrow[from=1-3, to=2-4]
	\arrow["{{(\src, \trg)}}", from=1-4, to=2-4]
	\arrow["{!}", from=2-2, to=1-1]
	\arrow["{{(0,1)}}"', from=2-2, to=1-2]
\end{tikzcd}
\]
Unfortunately, the map $(\src,\trg)$
is not as well-behaved as it was in the category of groupoids:
$\src : A^\two \to A$ is a fibration
whereas $\trg : A^\two \to A$ is an opfibration,
making $A^\two \to A \times A$
a two-sided fibration (or ``bifibration'' in \cite{street1974}).
% Two-sided fibrations are only classified bicategorically,
% adding a significant layer of complexity to the theory.
An approach to directed type theory using two-sided fibrations for
is explored in \cite{rivera2026}.

We consider the following alternative approach to the problem.
Recall the following definition of a twisted arrow category \cite{linton1965}.

\begin{definition}\label{def:twisted-arrow}
  For a category $A$,
  the \emph{twisted arrow category} $\Tw A$
  has morphisms of $A$ as its objects,
  and for two objects $(f : X \to Y)$ and $(g : W \to Z)$,
  a morphism $\phi : f \to g$ in $\Tw A$ consists of a pair of maps
  $\phi_0 : W \to X$ and $\phi_1 : Y \to Z$ such that
  $\phi_1 \circ f \circ \phi_0 = g$.
  \[ \begin{tikzcd}
    X & W \\
    Y & Z
    \arrow[""{name=0, anchor=center, inner sep=0}, "f"', from=1-1, to=2-1]
    \arrow["{\phi_0}"', from=1-2, to=1-1]
    \arrow[""{name=1, anchor=center, inner sep=0}, "g", from=1-2, to=2-2]
    \arrow["{\phi_1}"', from=2-1, to=2-2]
  \end{tikzcd}\]
  This is equipped with two endpoint projections
  $\src : \Tw{A} \to {A}^\op$,
  which takes an arrow to its source,
  and $\trg : \Tw{A} \to A$ which takes the target.
  There is also a functor
  $\rfl : \Core{A} \to \Tw{A}$ that takes
  an object $x$ in the groupoid core of $A$ to the identity on $x$
  and an isomorphism $f : x \to y$ to the twisted morphism
  $(f^{-1}, f) : \id_x \to \id_y$.
  \[\begin{tikzcd}
    x & {y} \\
    x & {y}
    \arrow["{=}"', from=1-1, to=2-1]
    \arrow["{f^{-1}}"', from=1-2, to=1-1]
    \arrow["{=}", from=1-2, to=2-2]
    \arrow["f"', from=2-1, to=2-2]
  \end{tikzcd}\]
\end{definition}

\begin{proposition}\label{prop:src-trg-discrete-opfibration}
  The pair of endpoint projections
  \[(\src,\trg) : \Tw{A} \to A^\op \times A\]
  is a discrete opfibration.
\end{proposition}
\begin{proof}
  Consider an object $(f : X \to Y)$ in $\Tw{A}$
  and a morphism $(\phi_0, \phi_1) : (X,W) \to (Y,Z)$ in $A^\op \times A$.
  The arrow
  $(\phi_0, \phi_1) : f \to \phi_1 \circ f \circ \phi_0$ in $\Tw{A}$
  is the unique lift of $(\phi_0, \phi_1)$.
  \[ \begin{tikzcd}
    X & W \\
    Y & Z
    \arrow[""{name=0, anchor=center, inner sep=0}, "f"', from=1-1, to=2-1]
    \arrow["{\phi_0}"', from=1-2, to=1-1]
    \arrow[""{name=1, anchor=center, inner sep=0}, "\phi_1 \circ f \circ \phi_0", from=1-2, to=2-2]
    \arrow["{\phi_1}"', from=2-1, to=2-2]
  \end{tikzcd}\]
\end{proof}

With the twisted arrow category $\Tw{A}$ replacing the arrow category $A^\two$,
we obtain an opfibration $\Tw{A} \to {A}^\op \times A$
factorising the \emph{twisted diagonal}
$(\inv, \inc) : \Core{A} \to {A}^\op \times A$,
which takes an object $x$ in the groupoid core of $A$ to
$(x,x)$ and an isomorphism $f : x \to y$
to a pair of maps $(f^{-1}, f) : (x,x) \to (y,y)$.
\[
\begin{tikzcd}
	{\Core A} & {\Tw{A}} \\
	& {{A}^\op \times A}
	\arrow["\rfl", from=1-1, to=1-2]
	\arrow["{(\inv,\inc)}"', from=1-1, to=2-2]
	\arrow["{{(\src,\trg)}}", from=1-2, to=2-2]
\end{tikzcd}
\]
We will call this the \emph{twisted pathobject factorisation}.


% In \Cref{def:twisted-arrow},
% we described the twisted arrow category on a category $A$,
% and the twisted pathobject factorisation of the map
% $(\inv,\inc) : \Core{A} \to {A}^\op \times A$.
% \[
% \begin{tikzcd}
% 	{\Core A} & {\Tw{A}} \\
% 	& {{A}^\op \times A}
% 	\arrow["\rfl", from=1-1, to=1-2]
% 	\arrow["{(\inv,\inc)}"', from=1-1, to=2-2]
% 	\arrow["{{(\src,\trg)}}", from=1-2, to=2-2]
% \end{tikzcd}
% \]

% Recall that a left adjoint $F$ in $F \dashv G : \CC \to \DD$
% is a ``left-adjoint-right-inverse'' or ``lari'' if
% the unit of the adjunction is the identity.

It will also be useful to define $\rTw{A}$,
the ``right twisted arrow category'' on $A$,
by the following pullback
\[ \begin{tikzcd}[column sep = large]
	{\rTw{A}} & {\Tw{A}} \\
	{(\Core{A}) \times A} & {{A}^\op \times A}
	\arrow[from=1-1, to=1-2]
	\arrow["{(\src,\trg)}"', from=1-1, to=2-1]
	\arrow["\lrcorner"{anchor=center, pos=0.125}, draw=none, from=1-1, to=2-2]
	\arrow["{(\src,\trg)}", from=1-2, to=2-2]
	\arrow["{\inv \times A}"', from=2-1, to=2-2]
\end{tikzcd} \]
Explicitly,
$\rTw{A}$ is the wide subcategory of the twisted arrow category $\Tw{A}$
such that morphisms $(\phi_0,\phi_1) : f \to g$ are in $\rTw{A}$
if and only if $\phi_0$ is invertible.
\[
\begin{tikzcd}
	x & w \\
	y & z
	\arrow["f"', from=1-1, to=2-1]
	\arrow["{\phi_0}"', from=1-2, to=1-1]
	\arrow["\sim", draw=none, from=1-2, to=1-1]
	\arrow["g", from=1-2, to=2-2]
	\arrow["{\phi_1}"', from=2-1, to=2-2]
\end{tikzcd}
\]
Then the twisted path factorisation can be pulled back, as follows,
to produce a \emph{right twisted pathobject factorisation}.
\begin{equation} \label{eq:global-right-twisted}
\begin{tikzcd}
	{\Core{A}} && \\
	& {\rTw{A}} & {\Tw{A}} \\
	& {(\Core{A}) \times A} & {{A}^\op \times A}
	\arrow["\rfl"{description}, dashed, from=1-1, to=2-2]
	\arrow["\rfl", bend left, from=1-1, to=2-3]
	\arrow["{\Delta := (\id, \inc)}"', bend right, from=1-1, to=3-2]
	\arrow[from=2-2, to=2-3]
	\arrow["{(\src,\trg)}"', from=2-2, to=3-2]
	\arrow["\lrcorner"{anchor=center, pos=0.125}, draw=none, from=2-2, to=3-3]
	\arrow["{(\src,\trg)}", from=2-3, to=3-3]
	\arrow["{\inv \times A}"', from=3-2, to=3-3]
\end{tikzcd}
\end{equation}
% Consider the functor
% $(\id_A, \inc) : \Core A \to \Core A \times A$ for a category $A$,
% where $\inc : \Core A \to A$ is the inclusion.
This factorisation is in fact equivalent to the canonical (lari, split opfibration)
factorisation given in \Cref{thm:lari-opfib-awfs}.
\[
\begin{tikzcd}
	& {\Delta / (\Core{A}) \times A} & \\
	{\Core{A}} && {(\Core{A}) \times A} \\
	& {\rTw{A}}
	\arrow["{R_\Delta}", from=1-2, to=2-3]
	\arrow["\simeq"{description, pos=0.3}, from=1-2, to=3-2]
	\arrow["{L_\Delta}", from=2-1, to=1-2]
	\arrow["{{\Delta}}"{description, pos=0.3}, from=2-1, to=2-3]
	\arrow["{{\text{lari}}}"{description}, from=2-1, to=3-2]
	\arrow["{{{\text{split opfibration}}}}"{description}, from=3-2, to=2-3]
\end{tikzcd}
\]

Furthermore, in this case, the second map
$\rTw{A} \to (\Core{A}) \times A$ is a discrete fibration,
as it is a pullback of one (\Cref{prop:src-trg-discrete-opfibration}).
Any lari functor is also initial (any left adjoint is initial),
so this coincides with the orthogonal
\emph{comprehensive factorisation}
of \cite{street1973}.

There is also a ``left twisted arrow category'' $\lTw{A}$,
defined similarly.
It is also an example of the (lari, split opfibration)
factorisation.
\[
\begin{tikzcd}
	{\lTw{A}} & {\Tw{A}} \\
	{{A}^\op \times \Core{A}} & {{A}^\op \times A}
	\arrow[from=1-1, to=1-2]
	\arrow["{(\src,\trg)}"', from=1-1, to=2-1]
	\arrow["\lrcorner"{anchor=center, pos=0.125}, draw=none, from=1-1, to=2-2]
	\arrow["{(\src,\trg)}", from=1-2, to=2-2]
	\arrow["{{A}^\op \times \inc}"', from=2-1, to=2-2]
\end{tikzcd}
\quad
\begin{tikzcd}
	{\Core{A}} && {{A}^\op \times \Core{A}} \\
	& {\lTw{A}}
	\arrow["{(\inv, \id)}", from=1-1, to=1-3]
	\arrow["{\text{lari}}"{description}, from=1-1, to=2-2]
	\arrow["{{\text{split opfibration}}}"{description}, from=2-2, to=1-3]
\end{tikzcd}
\]

We record these facts in the following proposition.
\begin{proposition}\label{prop:right-twisted-factorisation-equiv}
  The right twisted pathobject factorisation of the
  ``right diagonal'' \[\Core A \to (\Core{A}) \times A\] is equivalent
  to the (lari,split opfibration) factorisation (\Cref{thm:lari-opfib-awfs}),
  as well as the (initial, discrete opfibration) factorisation \cite{street1973}.
  % In particular, the left factor $\rfl : \Core A \to \rTw A$
  % is lari (hence initial),
  % and the right factor $(\src,\trg) : \rTw A \to \Core A \times A$
  % is a discrete opfibration (hence a split opfibration).
  The right adjoint of the lari functor $\rfl : \Core A \to \rTw A$
  is given by $\src : \rTw A \to \Core A$.

  Similarly, the left twisted pathobject factorisation of
  the map \[\Core A \to A^\op \times \Core A\] coincides with
  the (lari,split opfibration) factorisation,
  as well as the comprehensive factorisation.
  The right adjoint of the lari functor $\rfl : \Core A \to \lTw A$
  is given by $\trg : \lTw A \to \Core A$.
\end{proposition}

\section*{Relativised constructions} 
%  \label{sec:relativised-constructions}
Since we want to build a type-theoretic analysis of $\Cat$,
we are interested not just in constructions on objects $A$ of $\Cat$,
but on ``relativised objects'', i.e. (small) split opfibrations $A \to X$.
The analogue to this in the LCC setting would be working in a slice over $X$.
We will regard (covariant) type families over $X$ as split opfibrations over $X$.
Instead of considering the slice $\Cat / X$,
we will therefore be taking the category $\opfibcart(X)$ of split opfibrations over $X$
(for all type families),
or the category $\smlopfibcart(X)$ of small split opfibrations over $X$
(for small type families).
The lemmas in this section verify that the various constructions we
have considered so far can indeed be relativised in this way.

Taking the category of small split opfibrations and
opcartesian morphisms defines a contravariant pseudofunctor
from the 1-category of large categories
to the 2-category of extra-large categories,
\[\smlopfibcart(-) : \Cat^\op \to \tCat\]
where the action on morphisms $F : \DD \to \CC$ is given by
pullback along $F$.
It follows from \Cref{thm:opfibration-main} (or a small version thereof)
that $\smlopfibcart(-)$ is equivalent to the strict 2-functor
$[-,\cat] : \Cat^\op \to \tCat$,
using the fact that $\Cat$ is Cartesian closed\footnote{
  There are two subtle technicalities here.
  One is an issue of size.
  The category $\smlopfibcart(X)$ is extra-large,
  but is equivalent to $[X,\cat]$,
  which is just large.
  The second is an issue of dimension.
  Strictly speaking, the Cartesian closure provides a 1-functor
  $[-,\cat] : \Cat^\op \to \Cat$,
  which we compose with the inclusion $\Cat \subseteq \tCat$,
  which is a 2-functor.}.
\[\smlopfibcart(-) \simeq [-,\cat]\]
% \[
% \begin{tikzcd}
% 	{\Cat^\op} & \tCat \\
% 	& \Cat
% 	\arrow["{\smlopfibcart(-)}", from=1-1, to=1-2]
% 	\arrow["{[-,\cat]}"', from=1-1, to=2-2]
% 	\arrow["\simeq"', from=2-2, to=1-2]
% \end{tikzcd}
% \]
The twisted arrow structure can then be relativised --
constructed uniformly and locally over every object --
with $\smlopfibcart(X)$ playing the role of
a slice category $\CC / X$ in an LCC.
We outline this construction below.

Recall that there are endofunctors on $\cat$
that respectively take a category to its opposite,
a category to its core,
and a category to its twisted arrow category,
which are related by natural transformations,
variously taking source, target, and so on.
% (described in \Cref{sec:twisted-path-factorisation}).
Let us denote these functors and natural transformations as follows
(we add the overline annotation to reserve the
unannotated version for later).
\[
\overline{\vOp{}},\, \overline{\vCore{}}, \, \overline{\vTw{}}
: \cat \to \cat
\quad \quad 
\begin{tikzcd}
  & \overline{\vOp{}} \\
  \overline{\vCore{}} & \overline{\vTw{}} \\
  & {\id_{\cat}}
  \arrow["\overline\inv", from=2-1, to=1-2]
  \arrow["\overline\rfl"{description}, from=2-1, to=2-2]
  \arrow["\overline\inc"', from=2-1, to=3-2]
  \arrow["\overline\src"', from=2-2, to=1-2]
  \arrow["\overline\trg", from=2-2, to=3-2]
\end{tikzcd}
\]

Post-composition with $\overline{\vOp{}}$, $\overline{\vCore{}}$ and $\overline{\vTw{}}$
define the following (strict) 2-transformations
on $[-,\cat] : \Cat^\op \to \Cat$.
\[[-,\overline{\vOp{}}],\, [-,\overline{\vCore{}}], \, [-,\overline{\vTw{}}] : [-,\cat] \to [-,\cat]\]
% defines a (strict) 2-transformation
% \[[-,\overline{\vOp{}}] : [-,\cat] \to [-,\cat]\]
% Similarly, the core functor and twisted arrow functor
These are related by the following modifications.
\[
\begin{tikzcd}[column sep=huge]
  & {[-,\overline{\vOp{}}]} \\
  {[-,\overline{\vCore{}}]} & {[-,\overline{\vTw{}}]} \\
  & {\id}
  \arrow["{[-,\overline\inv]}", from=2-1, to=1-2]
  \arrow["{[-,\overline\rfl]}"{description}, from=2-1, to=2-2]
  \arrow["{[-,\overline\inc]}"', from=2-1, to=3-2]
  \arrow["{[-,\overline\src]}"', from=2-2, to=1-2]
  \arrow["{[-,\overline\trg]}", from=2-2, to=3-2]
\end{tikzcd}
\]

\begin{definition}\label{def:relativised-constructions}
Using the equivalence $[-,\cat] \simeq \smlopfibcart(-)$
we obtain three pseudotransformations
\[{\vOp{}}, {\vCore{}}, {\vTw{}}
: \smlopfibcart(-) \to \smlopfibcart(-) \]
which are related by the following modifications.
\[
\begin{tikzcd}
  & {\vOp{}} \\
  {\vCore{}} & {\vTw{}} \\
  & {\id_{\smlopfibcart(-)}}
  \arrow["\inv", from=2-1, to=1-2]
  \arrow["\rfl"{description}, from=2-1, to=2-2]
  \arrow["\inc"', from=2-1, to=3-2]
  \arrow["\src"', from=2-2, to=1-2]
  \arrow["\trg", from=2-2, to=3-2]
\end{tikzcd}
\]
\end{definition}
For each $X$ in $\Cat$,
the operation $\vOp[X]{} : \smlopfibcart(X) \to \smlopfibcart(X)$
applies $\vOp{}$ to each fibre of a split opfibration $A \to X$.
We therefore call $\vOp[X]{A}$ the \emph{vertical opposite} of $A$ (over $X$).
Similarly, we have the \emph{vertical core} $\vCore[X]{A}$ and \emph{vertical twisted arrows} $\vTw[X]{A}$.
\begin{definition}
  We apply the vertical core to the universe (over $\cat$) to produce
  \[ \usim := \vCore[\cat]{\ulow} : \cat_\sim \to \cat \]
\end{definition}
 
\begin{proposition}\label{prop:local-src-trg-discrete-opfibration}
  For a split opfibration $A \to X$,
  the pathobject map
  \[(\src,\trg) : \vTw[X] A \to \vOp[X]{A} \times_X A\]
  in the relativised twisted pathobject factorisation over $X$
  is a discrete opfibration,
  which is in particular a split opfibration.  
  When $A \to X$ is a small, so is $(\src,\trg)$.
\end{proposition}
\begin{proof}
  From \Cref{prop:src-trg-discrete-opfibration},
  \[(\src, \trg) : \Tw{(A^* x)} \to (A^* x)^\op \times A^* x\]
  is a discrete opfibration, for each $x$ in $X$.
  By \Cref{lem:split-opfib-between-split-opfibs},
  it follows that
  \[(\src,\trg) : \vTw[X] A \to \vOp[X]{A} \times_X A\]
  is a discrete opfibration.
  
  Suppose $A \to X$ is a small split opfibration.
  Then $\vTw[X] A \to X$ is small and
  factors as
  \[\vTw[X] A \to \vOp[X] A \times_X A \to X\]
  It follows that $\vTw[X] A \to \vOp[X] A \times_X A$ is small.
\end{proof}

What follows is a relativised version of the factorisation given in
\Cref{thm:lari-opfib-awfs}.
An explicit construction of this factorisation
is given in \cite[Lemma 4.33.14]{stacks-project}.
We provide an alternative construction,
by using straightening and unstraightening (\Cref{thm:opfibration-main})
and the global AWFS (\Cref{thm:lari-opfib-awfs}).
\begin{theorem}\label{thm:relativised-AWFS}
  Let $(A \to X,l)$ and $(B \to X,l)$
  be small split opfibrations over a category $X$.
  Any morphism of split opfibrations $F : A \to B$ over $X$
  factors in $\opfibcart(X)$ as follows.
  \[ \begin{tikzcd}
    A && B \\
    & {K_X F}
    \arrow["F", from=1-1, to=1-3]
    \arrow["{L_X F}"', from=1-1, to=2-2]
    \arrow["{R_X F}"', from=2-2, to=1-3]
  \end{tikzcd} \]
  The factorisation $(L_X F, R_X F)$
  forms the basis of an AWFS on $\opfibcart(X)$.

  Furthermore, these factorisations are stable under pullback
  along maps $G : Y \to X$.
  In particular, $K_F$ is the component of a modification
  $K : \opfibcart(-) \to \opfibcart(-)$.
  \[
  \begin{tikzcd}
    {\opfibcart(X)} & {\opfibcart(X)} \\
    {\opfibcart(Y)} & {\opfibcart(Y)}
    \arrow["K", from=1-1, to=1-2]
    \arrow["{G^*}"', from=1-1, to=2-1]
    \arrow["{G^*}", from=1-2, to=2-2]
    \arrow["K"', from=2-1, to=2-2]
  \end{tikzcd}
  \]
\end{theorem}
\begin{proof}
  To construct the monad
  $R_X : \opfibcart(X)^\two \to \opfibcart(X)^\two$
  we can use straightening and unstraightening
  (\Cref{thm:opfibration-main}),
  and the unrelativised AWFS on $\Cat$
  (\Cref{thm:lari-opfib-awfs}):
  \[ \begin{tikzcd}
    {\opfibcart(X)^\two} & {[X,\Cat]^\two} & {[X,\Cat^\two]} \\
    {\opfibcart(X)^\two} & {[X,\Cat]^\two} & {[X,\Cat^\two]}
    \arrow["\simeq"{description}, draw=none, from=1-1, to=1-2]
    \arrow["{R_X}"', dashed, from=1-1, to=2-1]
    \arrow["\simeq"{description}, draw=none, from=1-2, to=1-3]
    \arrow["{R \circ (-)}", from=1-3, to=2-3]
    \arrow["\simeq"{description}, draw=none, from=2-2, to=2-1]
    \arrow["\simeq"{description}, draw=none, from=2-3, to=2-2]
  \end{tikzcd} \]
  The unit and multiplication for $R_X$
  can be constructed by whiskering
  with the unit and multiplication for $R : \Cat^\two \to \Cat^\two$.
  The construction of $L_X$ is dual.
  % Briefly: by reindexing,
  % $F : \two \to [X,\cat]$ corresponds to a functor
  % \[ \tilde{F} : X \to \cat^\two\]
  % The comonad $L : \cat^\two \to \cat^\two$ and monad
  % $R : \cat^\two \to \cat^\two$
  % can be composed with $\tilde{F}$ to give
  % \[L \circ \tilde{F} : X \to \cat^\two \quad \text{and} \quad 
  %   R \circ \tilde{F} : X \to \cat^\two \]
  % These correspond to natural tranformations $L_F : A \to K_F$
  % and $R_F : K_F \to B$, when reindexing back to $\two \to [X,\cat]$.
\end{proof}

\begin{theorem}\label{thm:relativised-algebras}
  The coalgebras for the comonad $L_X : \opfibcart(X) \to \opfibcart(X)$
  and the algebras for the monad $R_X : \opfibcart(X) \to \opfibcart(X)$
  can be computed pointwise:
  \[
    \coalg{L_X} \simeq [X,\coalg{L}] \simeq [X,\lari]
    \quad \quad
    \alg{R_X} \simeq [X,\alg{R}] \simeq [X,\opfibcart]
  \]
  It follows from \Cref{lem:split-opfib-between-split-opfibs}
  and \Cref{lem:lari-between-split-opfibs}
  that an $R_X$-algebra structure on $F : A \to B$ over $X$
  produces a split opfibration structure on $F : A \to B$,
  and an $L_X$-algebra on $F$ produces a lari adjunction $F \dashv R$,
  where $R$ is a morphism of split opfibrations over $X$.
\end{theorem}

  % The category of coalgebras for the comonad $L_X : \opfibcart(X) \to \opfibcart(X)$
  % is equivalent to the category of functors $[X,\coalg{L}]$,
  % and the category of algebras for the monad $R_X : \opfibcart(X) \to \opfibcart(X)$
  % is equivalent to the category of $[X,\alg{R}]$.
  % $[X,\opfibcart]$
One can also construct a relativised right twisted arrow structure
\[\begin{tikzcd}
	& {\vCore{}} \\
	{\vCore{}} & {\rtw{}} \\
	& {\id_{\opfib(-)}}
	\arrow["{=}", from=2-1, to=1-2]
	\arrow["\rfl"{description}, from=2-1, to=2-2]
	\arrow["\inc"', from=2-1, to=3-2]
	\arrow["\src"', from=2-2, to=1-2]
	\arrow["\trg", from=2-2, to=3-2]
\end{tikzcd}\]

\begin{proposition}\label{prop:rtw-equiv-K}
  For a small opfibration $A \to X$,
  $\rtw[X]{A}$ and $K_\delta$ are equivalent as categories over $X$,
  where
  \[\delta = (\id,\inc) : \vCore[X]{A} \to \vCore[X]{A} \times_X A\] is the
  (right) diagonal and
  \[\begin{tikzcd}
    {\vCore[X]{A}} & {K_\delta} & {\vCore[X]{A} \times_X A}
    \arrow["{L_\delta}", from=1-1, to=1-2]
    \arrow["{R_\delta}", from=1-2, to=1-3]
  \end{tikzcd} \]
  is the functorial factorisation of $\delta$ given by the AWFS of
  \Cref{thm:relativised-AWFS}.
  Furthermore, this equivalence commutes with the two factorisations
  \[ \begin{tikzcd}
    & {\rtw[X]{A}} & \\
    {\vCore[X]{A}} && {\vCore[X]{A} \times_X A} \\
    & {K_\delta}
    \arrow["{(\src,\trg)}", from=1-2, to=2-3]
    \arrow["\simeq"{description}, from=1-2, to=3-2]
    \arrow["\rfl", from=2-1, to=1-2]
    \arrow["{L_\delta}"', from=2-1, to=3-2]
    \arrow["{R_\delta}"', from=3-2, to=2-3]
  \end{tikzcd}\]
  By \Cref{thm:relativised-algebras}, $L_\delta$ is lari.
  It follows that $\rfl : \vCore[X]{A} \to \rtw[X]{A}$
  has the left-lifting property with respect to
  split opfibrations.
\end{proposition}
\begin{proof}
  Write $\tilde{A} : X \to \cat$ for the fibres functor corresponding to
  the small opfibration $A$ in $\smlopfib(X)$.
  Then $\rtw[X]{A}$ in $\smlopfib(X)$ corresponds to the functor
  $\rtw{} \circ \tilde{A}$,
  where $\rtw{} : \cat \to \cat$ is the small version of the
  right twisted arrow functor $\rTw{} : \Cat \to \Cat$
  defined in \Cref{eq:global-right-twisted}.
  The equivalence of $K_{(\id,\inc)} \simeq \rtw{X}$ for a category $X$
  extends to an equivalence of 2-functors $K \simeq \rtw{}$
  in the 2-category $\cat \to \cat$.
  Whiskering provides an equivalence of 2-functors
  $K \circ \tilde{A} \simeq \rtw{} \circ \tilde{A}$,
  which translates to the desired equivalence $K_\delta \simeq \rtw[X]{A}$.
\end{proof}